\section{The First Week: Lifting Up Our Heads}\label{sec:week-1}

Advent begins on Saturday evening at First Vespers of the First Sunday
of Advent (RG~n.~71). The calendar fixes this Sunday as the one falling
on 30~November or nearest to it (RG~n.~20), placing it between
27~November and 3~December. Whenever 30~November itself is Sunday, the
First Sunday of Advent takes precedence as a first-class Sunday over the
second-class feast of St.~Andrew the Apostle (RG~nn.~11, 15, 91), and
St.~Andrew is omitted that year because a first-class day admits only a
privileged commemoration (RG~nn.~109, 111~a).

First Vespers already sounds the note of Sunday's Mass. Its short
chapter proclaims Romans~13:11: \textit{Hora est jam nos de somno
surgere}, ``it is now the hour for us to rise from sleep.'' The choir
sings the Advent Vespers hymn \textit{Creator alme siderum}, ``loving
Creator of the stars''---rendered by St.~John Henry Newman as
``Creator of the starry pole, / Saviour of all who live''
(Appendix~\ref{app:prayers})---followed by the versicle from
Isaiah~45:8: \textit{Rorate caeli desuper, et nubes pluant justum},
``Drop down, ye heavens, from above, and let the skies pour down the
Righteous One,'' and the reply, \textit{Aperiatur terra, et germinet
Salvatorem}, ``Let the earth open, and let her bring forth the
Saviour.'' Over the \textit{Magnificat}, the antiphon \textit{Ecce nomen
Domini venit de longinquo} announces that the name of the Lord comes
from afar and his brightness fills the whole earth. That night at
Compline, the Church begins \textit{Alma Redemptoris Mater}, ``Maiden!
Mother of Him Who redeemed us'' (Appendix~\ref{app:prayers}), pairing it
throughout Advent with the Annunciation versicle \textit{Angelus Domini
nuntiavit Mariae} and the prayer \textit{Gratiam tuam}.

\subsection{The Ox at the Crib in the Night Office}

At Matins on the First Sunday of Advent, the Roman Breviary opens the
book of Isaiah from its first line, \textit{Incipit liber Isaiae
Prophetae}. Under the 1960 rubrics (RGBR~nn.~162~a, 220), Sunday Matins
in Advent consists of one nocturn of nine psalms and three lessons: two
from Isaiah, uniting the older second and third lessons into one,
followed by St.~Gregory the Great's homily on the Gospel. In the third
verse of that opening lesson, four weeks before Christmas night, the
prophet delivers God's charge against an ungrateful people (Is~1:3):

\begin{quote}
The ox knoweth his owner, and the ass his master's crib: but Israel
hath not known me, and my people hath not understood.
\end{quote}

In the Latin Bible and the Breviary, the word for the feeding trough is
\textit{praesepe}: \textit{Cognovit bos possessorem suum, et asinus
praesepe domini sui}. It is the same word St.~Luke's Latin Gospel uses
when Mary wraps her newborn Son in swaddling bands and lays him in the
manger, \textit{reclinavit eum in praesepio} (Luke~2:7).

Commenting on Bethlehem (\textit{Exposition of the Gospel of Luke}
II.42--43), St.~Ambrose of Milan warns his hearers not to confine
divinity within the littleness of an infant body, and brings Isaiah's
beasts straight into the stable: \textit{quia in pannis est uides, quia
in caelis est non uides. infantis audis uagitus, non audis bouis dominum
agnoscentis mugitus; agnouit enim bos possessorem suum et asinus
praesaepe domini sui}---you see that he is in swaddling bands, you do
not see that he is in heaven; you hear an infant's crying, you do not
hear the lowing of the ox recognizing its Lord, for the ox has known his
owner, and the ass his master's crib. In the next section (II.43),
Ambrose opens the Eucharistic meaning of that feeding trough:
\textit{hic est dominus, hoc praesaepe, quo nobis diuinum mysterium
reuelatum est, inrationabili gentes pecudum intra praesaepia more
uiuentes alimoniae sacrae ubertate pascendos}---this is the Lord, this
the manger by which the divine mystery has been revealed to us, that the
nations who lived like unreasoning beasts within the stalls should be
fed with the abundance of sacred nourishment. In the East, St.~Gregory
Nazianzen preaches the same call (\textit{Oration} 38.17):

\begin{quote}
Know as Isaiah bids thee, thine Owner, like the ox, and like the ass
thy Master's crib; if thou be one of those who are pure and lawful food,
and who chew the cud of the word and are fit for sacrifice. Or if thou
art one of those who are as yet unclean and uneatable and unfit for
sacrifice, and of the gentile portion, run with the Star, and bear thy
Gifts with the Magi, gold and frankincense and myrrh, as to a King, and
to God, and to One Who is dead for thee.
\end{quote}

\subsection{At St.~Mary Major: Lifting the Soul and Waking from Sleep}

The Roman station for the First Sunday of Advent is the Basilica of
St.~Mary Major (\textit{Statio ad S.\ Mariam maiorem}), where the
Church begins the season and to which she will return on Ember
Wednesday, on the Vigil of Christmas, and at the Midnight and Day Masses
of the Nativity.

Throughout Advent, \textit{Gloria in excelsis} falls silent at Mass and
\textit{Te Deum} is omitted at Matins while the Sunday \textit{Alleluia}
remains, matching the mingled mourning and joy medieval liturgists saw in
the season (Section~\ref{sec:history}). In the 1962 rite, the deacon and
subdeacon wear dalmatic and tunicle rather than the abolished folded
chasubles (RG~n.~137); the organ is forbidden except to sustain singing at
Mass and Vespers, on holy days, and on Gaudete Sunday (Sacred Congregation
of Rites, \textit{Instructio de musica sacra}, 3~September 1958,
nn.~81~a, 83); and the Mass ends with \textit{Ite, missa est}
(RGMR~n.~507).

The Introit opens \textit{Ad te levavi}, ``To thee, O Lord, have I
lifted up my soul'' (Ps~24~[25]:1--3), begging in its verse, ``Shew, O
Lord, thy ways to me, and teach me thy paths'' (v.~4). St.~Augustine
(\textit{Expositions on the Psalms} 24.1--2) hears Christ speaking here
in the person of the Church: ``with spiritual longing have I lift up the
soul, that was trodden down on the earth with carnal longings.'' The
same psalm returns in the Gradual (\textit{Universi, qui te exspectant},
Ps~24~[25]:3--4) and the Offertory (\textit{Ad te levavi}), while the
Sunday Alleluia (\textit{Ostende nobis}, Ps~84~[85]:8) and the Communion
(\textit{Dominus dabit benignitatem}, Ps~84~[85]:13, ``For the Lord will
give goodness: and our earth shall yield her fruit'') open Psalm~84~[85].

After the \textit{Kyrie}, the priest prays the Collect of the First Week:

\begin{massbox}{Collect of the First Sunday of Advent}
\textbf{Latin.} Excita, quaesumus, Domine, potentiam tuam, et veni: ut
ab imminentibus peccatorum nostrorum periculis, te mereamur protegente
eripi, te liberante salvari: Qui vivis et regnas cum Deo Patre in
unitate Spiritus Sancti, Deus: per omnia saecula saeculorum. R.~Amen.

\textbf{English.} Bestir, O Lord, Thy might, we pray Thee and come;
That, defended by Thee, we may deserve rescue from approaching dangers
brought on by our sins, and being set free by Thee, obtain our
salvation. Who livest.
\end{massbox}

Addressing Christ directly (\textit{Qui vivis}, ``Who livest''), the
Church asks him to rouse his own power and deliver us from the dangers
brought on by our own sins. In the Secret (\textit{Haec sacra nos}) and
Postcommunion (\textit{Suscipiamus, Domine}), she prays that these
mysteries may cleanse us to approach their Author and prepare us with
due reverence for the coming festival of our redemption.

The Epistle (Rom~13:11--14) rouses the congregation directly: ``And
that, knowing the season, that it is now the hour for us to rise from
sleep. For now our salvation is nearer than when we believed.'' St.~Paul
bids Christians cast off the works of darkness, put on the armour of
light, walk honestly as in the day---``not in rioting and drunkenness,
not in chambering and impurities, not in contention and envy''---closing
the liturgical reading with the command, ``put ye on the Lord Jesus
Christ.''

Preaching on this passage (\textit{Homily} 24 on Romans), St.~John
Chrysostom notes how easily early zeal grows slack, whereas Paul presses
us forward as the hour draws near: ``For the nearer the King may be at
hand, the more ought they to get themselves in readiness; the nearer the
prize is, the more wide awake ought they to be for the contest, since
even the racers do this, when they are upon the end of the course, and
towards the receiving of the prize, then they rouse themselves up the
more.'' Like householders who wake one another when the swallow twitters
before dawn so daylight does not catch them still in bed, we are called
to put on ``the Lord Himself for a garment, the King Himself,'' that he
may dwell within the soul and clothe us on every side.

That same chapter of Romans---extending through the closing clause of
verse~14---broke open one of the decisive conversions in Christian
history. In the garden at Milan in the summer of 386, torn by years of
delay, St.~Augustine heard the child's voice bidding him take and read,
and hurried back to the volume of St.~Paul (\textit{Confessions}
VIII.12.29):

\begin{quote}
So quickly I returned to the place where Alypius was sitting; for
there had I put down the volume of the apostles, when I rose thence. I
grasped, opened, and in silence read that paragraph on which my eyes
first fell,---``Not in rioting and drunkenness, not in chambering and
wantonness, not in strife and envying; but put ye on the Lord Jesus
Christ, and make not provision for the flesh, to fulfil the lusts
thereof.'' No further would I read, nor did I need; for instantly, as the
sentence ended,---by a light, as it were, of security infused into my
heart,---all the gloom of doubt vanished away.
\end{quote}

Preaching on Romans~13:11 (\textit{Parochial and Plain Sermons}, Vol.~I,
Sermon~5), St.~John Henry Newman turns the same call onto baptized
Christians living in quiet ease:

\begin{quote}
Many live altogether as though the day shone not on them, but the
shadows still endured; and far the greater part of them are but very
faintly sensible of the great truths preached around them. They see and
hear as people in a dream; they mix up the Holy Word of God with their
own idle imaginings; if startled for a moment, still they soon relapse
into slumber; they refuse to be awakened, and think their happiness
consists in continuing as they are.
\end{quote}

\subsection{When the World Shakes: The Gospel of the First Sunday}

The Gospel (Luke~21:25--33) opens at the end of time, with ``signs in
the sun and in the moon and in the stars; and upon the earth distress of
nations, by reason of the confusion of the roaring of the sea, and of
the waves,'' men withering away for fear as the powers of heaven are
moved and the Son of man comes in a cloud with great power and majesty.
Yet right inside that upheaval, Christ commands his disciples
(Luke~21:28): ``But when these things begin to come to pass, look up and
lift up your heads, because your redemption is at hand.''

Pope St.~Gregory the Great preached his first homily on the Gospels on
this passage, and the Roman Breviary reads its opening at Sunday Matins
(\textit{Homily} 1.1): ``Our Lord and Saviour wisheth to find us ready
at His second coming. Therefore He telleth us what will be the evils of
the world as it groweth old, that He may wean our hearts from worldly
affections. Here we read what great convulsions will go before the end,
that, if we will not fear God in our prosperity, we may at least be
scourged into fearing His judgment when it is at hand.'' Then, on
verse~28, Gregory explains why Christians lift their heads when the
world shakes (\textit{Homily} 1.3, in the \textit{Catena Aurea} on
Luke~21:28):

\begin{quote}
Having in what has gone before spoken against the reprobate, He now
turns His words to the consolation of the elect; for it is added, When
these things begin to be, look up, and lift up your heads, for your
redemption draweth nigh; as if He says, When the buffetings of the world
multiply, lift up your heads, that is, rejoice your hearts, for when the
world closes whose friends ye are not, the redemption is near which ye
seek. For in holy Scripture the head is often put for the mind, for as
the members are ruled by the head, so are the thoughts regulated by the
mind. To lift up our heads then, is to raise up our minds to the joys of
the heavenly country.
\end{quote}

\subsection{Keeping the First Week Day by Day}

From Monday through Saturday, ordinary weekdays of Advent have no proper
Mass of their own. Whenever no saint's feast intervenes, the priest
repeats Sunday's Mass, keeping the Gradual and omitting the
\textit{Alleluia} and its verse with no replacement Tract
(RGMR~n.~299). Through 16~December, ordinary Advent weekdays rank as
third-class ferias and yield to feasts such as St.~Andrew
(30~November), St.~Bibiana (2~December), and St.~Francis Xavier
(3~December), while always receiving a privileged commemoration
(RG~nn.~25~b, 109~e).

Meanwhile, Matins advances daily through Isaiah~1--7. Monday continues
chapter~1 with the call to repentance and cleansing, ``Wash yourselves, be
clean'' (Is~1:16--18; Section~\ref{sec:penance-law}). Tuesday reads of
the mountain of the Lord to which all nations stream (Is~2); Wednesday
and Thursday trace judgment and cleansing in Sion (Is~3--4); Friday hears
Isaiah's vision of the Lord in the temple and the seraph's live coal
(Is~6); and Saturday opens chapter~7, where King Achaz receives the sign
of the Virgin and Emmanuel (Is~7:14).

Keep this week at home in step with Mass and Office (following
Appendix~\ref{app:daily-course} for each day's readings and prayer).
Read Isaiah~1--7 alongside Romans~13:11--14 and Luke~21:25--33, and
spend a quiet quarter-hour with St.~Gregory's first homily on the
Gospels or Book~VIII of St.~Augustine's \textit{Confessions}. Pray the
\textit{Excita} Collect each morning until its words live in your
memory, and close each evening with \textit{Alma Redemptoris Mater}
(Appendix~\ref{app:prayers}). Put St.~Paul's wakefulness into concrete
practice in the household: rise promptly without lingering, cut away one
evening habit of excess or distraction that dulls prayer, and examine
which earthly attachment would make you dread the Lord's coming rather
than lift up your head to greet him.
